% Hafta 10 — Kuantum sonrası: neyi değiştiriyor, neyi değiştirmiyor?
% Derleme: py -3.12 tools/latex/derle.py docs/week-10/assets/h10-15-kuantum-sonrasi.tex
\documentclass[tikz,border=8pt]{standalone}
\usepackage{cen429-cizim}
\begin{document}
\begin{tikzpicture}[node distance=10mm and 10mm]
  \node[baslik] (b) at (0,0) {Kuantum sonrası: neyi değiştiriyor, neyi değiştirmiyor?};
  \node[altbaslik, text width=170mm] at (b.south west) {NIST seçimi: ML-KEM (Kyber) ve ML-DSA (Dilithium)};

  \node[kotukutu=78mm, below=10mm of b.south west, anchor=north west] (sol) {\kb{TEHDİT ALTINDA}\\RSA, DH, ECDH, ECDSA\\``şimdi topla, sonra çöz''\\[2mm]bugünkü trafik ileride çözülebilir\\geçiş bugün planlanır};
  \node[iyikutu=78mm, right=10mm of sol] (sag) {\kb{BÜYÜK ÖLÇÜDE GÜVENLİ}\\AES-256 (anahtar boyu yeter)\\SHA-384 / SHA-3\\[2mm]simetrik dünya ayakta\\yalnız anahtar boyu artar};
  \draw[draw=cizgi, dashed] ($(sol.north east)!0.5!(sag.north west)$) -- ($(sol.south east)!0.5!(sag.south west)$);

  \node[dip, text width=170mm, below=10mm of sag.south -| sol, anchor=north west] {%
    Uzun ömürlü sırlarınız varsa hibrit anahtar değişimine geçiş bugünün kararıdır.};
\end{tikzpicture}
\end{document}
